\documentclass[12pt]{amsart}
%prepared in AMSLaTeX, under LaTeX2e
\addtolength{\oddsidemargin}{-.65in} 
\addtolength{\evensidemargin}{-.65in}
\addtolength{\topmargin}{-.55in}
\addtolength{\textwidth}{1.3in}
\addtolength{\textheight}{.9in}

\renewcommand{\baselinestretch}{1.05}

\usepackage{verbatim} % for "comment" environment

\usepackage{palatino}

\newtheorem*{thm}{Theorem}
\newtheorem*{defn}{Definition}
\newtheorem*{example}{Example}
\newtheorem*{problem}{Problem}
\newtheorem*{remark}{Remark}

\usepackage{fancyvrb,xspace,dsfont}

\usepackage[final]{graphicx}

% macros
\usepackage{amssymb}
\usepackage[dvipsnames]{xcolor}

\usepackage{hyperref}
\hypersetup{pdfauthor={Ed Bueler},
            pdfcreator={pdflatex},
            colorlinks=true,
            citecolor=blue,
            linkcolor=red,
            urlcolor=blue,
            }

\newcommand{\br}{\mathbf{r}}
\newcommand{\bv}{\mathbf{v}}
\newcommand{\bx}{\mathbf{x}}
\newcommand{\by}{\mathbf{y}}

\newcommand{\cB}{\mathcal{B}}
\newcommand{\cD}{\mathcal{D}}
\newcommand{\cF}{\mathcal{F}}
\newcommand{\cH}{\mathcal{H}}
\newcommand{\cL}{\mathcal{L}}
\newcommand{\cV}{\mathcal{V}}
\newcommand{\cW}{\mathcal{W}}

\newcommand{\CC}{\mathbb{C}}
\newcommand{\NN}{\mathbb{N}}
\newcommand{\RR}{\mathbb{R}}
\newcommand{\ZZ}{\mathbb{Z}}

\renewcommand{\Im}{\mathrm{Im}}
\renewcommand{\Re}{\mathrm{Re}}

\newcommand{\eps}{\epsilon}
\newcommand{\grad}{\nabla}
\newcommand{\lam}{\lambda}
\newcommand{\lap}{\triangle}

\newcommand{\ip}[2]{\ensuremath{\left<#1,#2\right>}}

\newcommand{\image}{\operatorname{im}}
\newcommand{\onull}{\operatorname{null}}
\newcommand{\rank}{\operatorname{rank}}
\newcommand{\range}{\operatorname{range}}
\newcommand{\trace}{\operatorname{tr}}
\newcommand{\Span}{\operatorname{span}}

\newcommand{\prob}[1]{\bigskip\noindent\textbf{#1.}\quad }

\newcommand{\pts}[1]{(\emph{#1 pts}) }
\newcommand{\epart}[1]{\medskip\noindent\textbf{(#1)}\quad }
\newcommand{\ppart}[1]{\,\textbf{(#1)}\quad }

\newcommand{\ds}{\displaystyle}

\newcommand{\nex}{\medskip\noindent}


\begin{document}
\scriptsize \noindent Math 617 Functional Analysis (Bueler) \hfill \emph{assigned February 2026}
\normalsize\medskip

\Large\centerline{\textbf{Assignment 3}}
\large
\medskip

\centerline{\textbf{Due XX February at the beginning of class}}
\medskip
\normalsize

\thispagestyle{empty}

\bigskip
\noindent This Assignment is based on Chapters 2 and 4 of our textbook,\footnote{K.~Saxe,~\emph{Beginning Functional Analysis}, Springer 2010.} and on the lectures.  When you do calculations from Chapter 4 you will use the $L^2$ spaces, but knowledge of the Lebesgue integral will not be necessary; just regard integrals as Riemann integrals if you wish.


\newcommand{\augment}[1]{\quad\,\strut \begin{minipage}[t]{0.7\textwidth} \emph{#1}\end{minipage}}
\newcommand{\augpart}[1]{\emph{\textbf{(#1)}}}

\bigskip
\noindent \textsc{Do the following} Exercises from Chapter 2 (see pages 29--31):

\begin{itemize}
\item \textbf{Exercise 2.2.3}
\item \textbf{Exercise 2.3.4}
\item \textbf{Exercise 2.3.7}
\end{itemize}

\bigskip
\noindent \textsc{Do the following} Exercises from Chapter 4 (see pages x--y):

\begin{itemize}
\item \textbf{Exercise 4.1.1} \augment{Note that ``converges in mean'' is a synonym for ``converges in $\|\cdot\|_2$''.}
\item \textbf{Exercise 4.1.3} \augment{For part \augpart{b} please justify the use of a theorem.  For part \augpart{c}, please give the correct answer even if you cannot prove it.  (Perhaps graph some partial sums?  Or look up theorems from other sources.)}
\item \textbf{Exercise 4.1.4} \augment{Part \augpart{c} needs the fact that $[-\pi,\pi]$ is of finite length.}
\item \textbf{Exercise 4.2.2} \augment{You may assume that ``$L^2$'' refers to $L^2([a,b])$ for $a,b\in\RR$, but the statement holds generally, for any $L^2(X,\mu)$.  Also, the hypotheses \emph{do} imply that $f\in L^2$.  (Why?)}
\item \textbf{Exercise 4.2.3}
\item \textbf{Exercise 4.2.5} \augment{This calculation was mostly done in the week 4 slides.  See the \href{https://bueler.github.io/fa/daily.html}{Daily Log tab} for the blank slides PDF.}
\end{itemize}


\medskip
\noindent \textsc{Do the following additional problems.}

\prob{P4}  \emph{Chapter 2 defines \emph{equicontinuous} and proves the Ascoli-Arzel\`a Theorem.  However, it gives no related exercises!  Here is one.}

\epart{a}  Consider $V=C([0,1])$ with the supremum norm $\|\cdot\|_\infty$.  For $\eps>0$ and $\alpha \in (0,1)$ let
	$$f_{\alpha,\eps}(x) = \begin{cases} 0, & 0 \le x \le \alpha \\
	                                     (x-\alpha)/\eps, & \alpha \le x \le \alpha + \eps \\
	                                     1, & \alpha + \eps \le x \le 1 \end{cases}$$
	                             

\medskip \noindent yy

\end{document}
